\section{Implementation}

The platform is built entirely from version-controlled Infrastructure-as-Code,
in six stages: provision the VMs, bootstrap the cluster, install the secondary
CNI, deploy the 5G core, attach a simulated RAN, and add observability. The
first three build the \emph{NFVI}; the last three put a real 5G core on it.

\subsection{Infrastructure as Code (Vagrant)}

A single \texttt{Vagrantfile} declares the three VMs, each with two internal
networks (management and data). Critically, promiscuous mode is enabled on the
data NIC so macvlan works across nodes. A common \texttt{bootstrap.sh} disables
swap, loads the \texttt{br\_netfilter} and \texttt{overlay} kernel modules,
applies the Kubernetes sysctls, and sets a reliable DNS resolver.

\begin{lstlisting}[language=Ruby, caption=Per-node network + promiscuous mode (excerpt)]
config.vm.provider "virtualbox" do |vb|
  vb.cpus = 2 ; vb.memory = 4096
  # macvlan cross-node trên data NIC (nic3 = eth2) cần nhận frame của MAC con
  vb.customize ["modifyvm", :id, "--nicpromisc3", "allow-all"]
end
node.vm.network "private_network", ip: "10.10.10.10", virtualbox__intnet: "mgmt-net"  # eth1
node.vm.network "private_network", ip: "10.30.0.10",  virtualbox__intnet: "data-net"  # eth2
\end{lstlisting}

\subsection{Cluster provisioning (Ansible / k3s)}

Two idempotent playbooks install the cluster. \texttt{k3s-master.yml} installs
the k3s server on node0, pinning the node IP and Flannel interface to the
management network so control-plane signalling never crosses the data plane.
Traefik and the bundled load balancer are disabled to save RAM; metrics-server
is kept because the HPA needs it.

\begin{lstlisting}[language=bash, caption=k3s server install flags]
INSTALL_K3S_EXEC="server \
  --node-ip 10.10.10.10 --advertise-address 10.10.10.10 \
  --flannel-iface eth1 --disable traefik --disable servicelb \
  --write-kubeconfig-mode 644 --tls-san 10.10.10.10"
\end{lstlisting}

\texttt{k3s-worker.yml} joins node1 and node2 as agents on the management
interface. After this stage \texttt{kubectl get nodes} shows three
\texttt{Ready} nodes --- the NFVI compute is live.

\subsection{Secondary CNI (Multus + Whereabouts)}

\texttt{cni.yml} installs the thin Multus daemonset and Whereabouts IPAM. Two
k3s-specific adaptations are required: (i)~the \texttt{hostPath} values are
rewritten to k3s's CNI directories
(\texttt{/var/lib/rancher/k3s/\ldots}) without touching the in-container mount
paths; and (ii)~\texttt{macvlan} (plus \texttt{static}, \texttt{ipvlan}) are
symlinked onto the k3s multicall CNI binary on every node, since k3s ships
symlinks for only a subset of plugins. This is the infrastructure that will let
the UPF obtain its user-plane interfaces.

\subsection{Deploying the 5G core (Open5GS)}

The 5G core is deployed as CNFs using the community
\texttt{towards5gs-helm} charts \cite{towards5gs}, which package Open5GS for
Kubernetes with Multus. The \texttt{open5gs.yml} playbook installs the core
with a values file tuned for this testbed: the N2/N3/N4/N6 networks are bound to
\texttt{eth2} (Section~4.2), every NF is given small CPU/memory requests, and
MongoDB persistence is disabled.

\begin{lstlisting}[language=yaml, caption=open5gs-values.yaml --- binding 5G interfaces to the data NIC (excerpt)]
global:
  n3network: { masterIf: eth2, subnetIP: 10.30.30.0, cidr: 24 }   # GTP-U
  n6network: { masterIf: eth2, subnetIP: 10.30.60.0, cidr: 24 }   # to DN
upf:
  n3if: { ipAddress: 10.30.30.10 }
  n4if: { ipAddress: 10.30.40.10 }
  n6if: { ipAddress: 10.30.60.10 }
  config: { upf: { pdn: [ { addr: 10.45.0.1/16, dnn: internet } ] } }  # IP pool cấp cho UE
\end{lstlisting}

After \texttt{helm install}, \texttt{kubectl get pods -n open5gs} shows the
full core --- NRF, SCP, AMF, SMF, UPF, AUSF, UDM, UDR, PCF, NSSF, BSF, and
MongoDB --- each as its own Deployment. The NFs register with the NRF over SBI;
the UPF comes up with its \texttt{n3}/\texttt{n4}/\texttt{n6} interfaces.

\subsection{Simulated RAN and subscriber provisioning (UERANSIM)}

A real UE requires a SIM record (IMSI, key, OPc) in the subscriber database.
A Kubernetes \texttt{Job} runs \texttt{open5gs-dbctl} to insert one test
subscriber (IMSI \texttt{001010000000001}) into MongoDB. \textbf{UERANSIM} is
then deployed (the \texttt{ueransim} chart) as a gNodeB plus a UE whose
identity matches that subscriber. The gNB attaches to the AMF over N2; the UE
performs registration and requests a PDU session.

\begin{lstlisting}[language=bash, caption=Order of deployment (Ansible open5gs.yml)]
helm install open5gs charts/open5gs   -n open5gs -f helm/open5gs-values.yaml
kubectl apply -f manifests/open5gs/subscriber-job.yaml   # nap SIM vao MongoDB
helm install ueransim charts/ueransim -n open5gs -f helm/ueransim-values.yaml
\end{lstlisting}

\subsection{Observability and access}

The \texttt{kube-prometheus-stack} chart is installed into the
\texttt{monitoring} namespace with a RAM-tuned values file (Alertmanager off,
persistence off, retention 2~days, Grafana on NodePort~30030). Because the
cluster networks are VirtualBox internal networks with no route from the host,
Grafana is reached through an SSH tunnel:

\begin{lstlisting}[language=bash, caption=SSH tunnel for Grafana access (PowerShell on Windows host)]
vagrant ssh node0 -- -N -L 3000:localhost:30030
# then browse http://localhost:3000  (admin / admin)
\end{lstlisting}
